\documentclass[../Main.tex]{subfiles}

\begin{document}

\section{Einleitung}

\subsection{Motivation und Relevanz}

Diese Maturitätsarbeit entstand aus dem Wunsch, Spiele nicht nur zu spielen, sondern sie genauer zu verstehen und selbst ein solches Spielerlebnis zu erschaffen.

Das Hack-and-Slash-Genre stellt aus unserer Sicht besonders hohe Anforderungen an das Zusammenspiel einzelner Spielsysteme, da Spielerbewegung, Kampfsystem, Gegnerverhalten und Leveldesign in Echtzeit ineinandergreifen. Im Zentrum unseres Interesses stehen das Gegnerverhalten und die Gegnerplatzierung. Wir haben diese beiden Bereiche gewählt, weil das Genre auf dem Kampfgeschehen aufbaut und die Gegner deshalb unserer Einschätzung nach für das Spielerlebnis besonders tragend sind.

Zu Videospielen gibt es umfangreiche Forschung und Fachliteratur. Auch im Rahmen dieser Maturitätsarbeit haben wir auf zahlreiche Quellen zurückgegriffen, etwa auf das MDA-Framework\autocite{hunicke2004mda}, auf Vorträge der Game Developers Conference (GDC) oder auf Methoden der Games User Research\autocite{drachen2018gamesuserresearch}. Diese werden in Kapitel~\ref{sec:Grundlagen} erläutert. Im Unterschied zu unserer Arbeit behandeln diese Quellen jedoch meist Spiele im Allgemeinen und nicht spezifisch das Hack-and-Slash-Genre.

In unserer Wahrnehmung erschienen Genrevertreter wie \game{Ninja Gaiden Sigma Plus} oder \game{NieR: Automata} in einer Phase, in der das Genre vergleichsweise häufig neue Titel hervorbrachte, während heute nur noch wenige grosse Titel erscheinen. Dadurch gibt es einerseits genügend Referenzmaterial, um zentrale Genremerkmale herauszuarbeiten. Andererseits sind wir bei der Umsetzung nicht an aktuelle Genrekonventionen gebunden, was uns zusätzlichen Freiraum für eigene, moderne Interpretationen klassischer Mechaniken lässt.

Unser Prototyp richtet sich in erster Linie an Fans der oben genannten Spiele. Da das Hack-and-Slash-Genre ein Untergenre der Action-Spiele ist, zählen wir auch Fans von Action-Spielen allgemein zu unserer Zielgruppe.

\begin{itemize}[leftmargin=1.5cm]
    \item \textcolor{red}{\textbf{Forschungsstand:} Optional noch die Leveldesign-Literatur (Totten, Rogers) und Game Feel (Swink) nennen, da sie für Kapitel~\ref{sec:Leveldesign} und~\ref{sec:Umsetzung} zentral sind. Fussnoten bei MDA und Drachen habe ich gesetzt; bei den GDC-Talks reicht der Verweis auf Kapitel~\ref{sec:Grundlagen}.}
    \item \textcolor{red}{\game{Ninja Gaiden Sigma Plus} und \game{NieR: Automata} in der Ludografie aufführen und beim ersten Nennen mit Fussnote zitieren.}
    \item \textcolor{red}{Optional für einen spannenden Einstieg: mit einer konkreten Spielszene beginnen (z.\,B. ein Gegner in \game{Ninja Gaiden}, der dich flankiert, während du einen anderen bekämpfst).}
\end{itemize}

\subsection{Erkenntnisfrage} \label{sec:Erkenntnisfrage}

Im Zentrum unserer Arbeit steht folgende Erkenntnisfrage:

\vspace{\baselineskip}

\begin{quote}
    \textit{Wie müssen Gegnerverhalten und Gegnerplatzierung in einem 3D-Hack-and-Slash-Prototyp konzipiert sein, damit ein zusammenhängender Kampfabschnitt mit nachvollziehbarem Spielfluss, angemessener Schwierigkeit und befriedigendem Spielgefühl entsteht?}
\end{quote}

\vspace{\baselineskip}

Diese Frage setzt sich aus zwei inhaltlichen Komponenten zusammen: Zum einen geht es um das Gegnerverhalten und das Kampfgefühl, zum anderen um die Gegnerplatzierung und das Leveldesign. Beide Komponenten werden anhand von drei Qualitäten beurteilt: Spielfluss, Schwierigkeit und Spielgefühl. Wie wir diese Begriffe verstehen, erläutern wir in Kapitel~\ref{sec:Grundlagen}.

\begin{itemize}[leftmargin=1.5cm]
    \item \textcolor{red}{\textbf{Vertrag:} Die Erkenntnisfrage weicht vom Vertrag ab (dort: \enquote{Wie können \ldots so umgesetzt werden, dass ein Kampfabschnitt von Testspieler:innen als flüssig und herausfordernd erlebt wird?}). Mit dem Betreuer klären oder angleichen.}
    \item \textcolor{red}{\textbf{Begriffe:} Der Verweis ist jetzt da, aber Kapitel~\ref{sec:Grundlagen} definiert zurzeit nur das Spielgefühl (Game Feel nach Swink, Abschnitt~2.2). Spielfluss und Schwierigkeit müssen dort noch definiert werden, sonst geht der Verweis ins Leere.}
    \item \textcolor{red}{\textbf{Kampfgefühl vs.\ Spielgefühl:} Wenn das zwei verschiedene Dinge sind, kurz den Unterschied erklären. Wenn nicht, ein Wort verwenden (dann auch in der Einleitung von Kapitel~\ref{sec:Umsetzung} anpassen).}
\end{itemize}

\subsection{Zielsetzung und Aufbau der Arbeit} \label{sec:Aufbau}

\textcolor{red}{Am Ende den Vertrag, das Spiel und den Text unten Abgleichen}

Ziel dieser Arbeit ist es, unsere Erkenntnisfrage anhand eines selbst entwickelten, spielbaren 3D-Hack-and-Slash-Prototyps zu beantworten. Auf dem Weg dorthin haben wir uns drei Teilziele gesetzt:

\begin{enumerate}
    \item die Analyse zweier Referenzspiele,
    \item die Entwicklung eines spielbaren Prototyps mit einem zusammenhängenden Kampfabschnitt,
    \item die Evaluation des Prototyps durch Playtests.
\end{enumerate}

Dieser Begleitbericht gehört zu unserem Prototyp.\footnote{\textcolor{red}{[Download-Link bzw.\ Bezugsquelle des Prototyps einfügen]}} Ergänzend dazu gibt es ein separates Dokument, das die technische Entwicklung in Godot dokumentiert.

Um diese Ziele zu erreichen, sind wir wie folgt vorgegangen: Zuerst haben wir ausgehend von der Erkenntnisfrage zu den Themen Gegnerverhalten (Kapitel~\ref{sec:Referenzspielanalyse}) und Gegnerplatzierung (Kapitel~\ref{sec:Leveldesign}) recherchiert. Um das Gegnerverhalten zu untersuchen, haben wir Genrevertreter selbst gespielt und unsere Beobachtungen mithilfe eines eigenen Analyseframeworks (Abschnitt~\ref{sec:Analysefragebogen}) festgehalten. Für die Gegnerplatzierung haben wir Fachliteratur zum Leveldesign sowie Vorträge von Fachleuten ausgewertet. Daraus haben wir acht Designvorschläge abgeleitet, je vier zum Gegnerverhalten und zur Gegnerplatzierung, die nach unserer Einschätzung besonders viel zur Beantwortung der Erkenntnisfrage beitragen. Diese Vorschläge haben wir anschliessend konkret in unserem Prototyp umgesetzt (Kapitel~\ref{sec:Umsetzung}).

In Playtests mit fünf Testpersonen haben wir dann versucht, die Designvorschläge zu evaluieren (Kapitel~\ref{sec:Evaluation}). Zum Schluss haben wir aus der Evaluation, den Playtests und unseren Recherchen das weitere Vorgehen abgeleitet und Faktoren zusammengetragen, die nach unserer Einschätzung das Gegnerverhalten und die Gegnerplatzierung ebenfalls beeinflussen und die wir als Nächstes untersuchen würden (Kapitel~\ref{sec:Schlussfolgerungen}).

\begin{itemize}[leftmargin=1.5cm]
    \item \textcolor{red}{\textbf{Teilziele überprüfbar machen:} Mit etwas mehr Detail lässt sich in Kapitel~\ref{sec:Schlussfolgerungen} klar sagen, ob sie erreicht wurden, z.\,B. (1) \enquote{mithilfe eines begründeten Analyseframeworks}, (2) \enquote{mit mindestens zwei unterscheidbaren Gegnertypen}, (3) \enquote{mit rund fünf Testpersonen}. So steht es auch im Vertrag.}
    \item \textcolor{red}{\textbf{Produkt} in 1--2 Sätzen beim Absatz \enquote{Dieser Begleitbericht \ldots}: Godot, Third-Person, ein Kampfabschnitt aus zwölf Räumen, fünf Gegnertypen; für Details auf Abschnitt~\ref{sec:Prototyp} und das README verweisen.}
    \item \textcolor{red}{\textbf{Technische Dokumentation:} Titel des Dokuments nennen, damit der Leser es findet.}
    \item \textcolor{red}{\textbf{Abgrenzung} (1 Satz): Prototyp statt fertiges Spiel; qualitative Playtests statt Statistik.}
    \item \textcolor{red}{\textbf{Aufbau unvollständig:} Kapitel~\ref{sec:Grundlagen} (Grundlagen und Methodik), Kapitel~\ref{sec:Reflexion} (Reflexion) und der Anhang mit den Playtest-Protokollen (Abschnitt~\ref{sec:PlaytestProtokolle}) fehlen noch. Je ein Satz genügt.}
    \item \textcolor{red}{Kapitel~\ref{sec:Referenzspielanalyse} behandelt auch die Gegnerplatzierung (Fragen A3, D3). Ggf. \enquote{vor allem} ergänzen.}
\end{itemize}

\newpage

\end{document}
